% Generated draft: requires mathematical and source review.
\documentclass[11pt]{article}
\usepackage[margin=25mm]{geometry}
\usepackage{fontspec}
\setmainfont{Latin Modern Roman}
\usepackage{amsmath,amssymb,mathtools,unicode-math}
\setmathfont{Latin Modern Math}
\usepackage{enumitem,needspace,xurl,hyperref}
\setlength{\parindent}{0pt}
\setlength{\parskip}{5pt}
\setlength{\emergencystretch}{3em}
\newcommand{\N}{\mathbb{N}}
\newcommand{\Z}{\mathbb{Z}}
\newcommand{\Q}{\mathbb{Q}}
\newcommand{\R}{\mathbb{R}}
\newcommand{\C}{\mathbb{C}}
\newcommand{\F}{\mathbb{F}}
\newcommand{\A}{\mathbb{A}}
\newcommand{\PP}{\mathbb{P}}
\newcommand{\E}{\mathbb{E}}
\newcommand{\eps}{\varepsilon}
\newcommand{\dd}{\,\mathrm{d}}
\begin{document}
\section*{UnsolvedMath ID: 30005298}
\textbf{Source problem code:} catalogue entry
\subsection{Definitions and mathematical statement}
Write \(Z=\{p_1,\ldots,p_r\}\subset\mathbb P^n\) for a set of \(r\) distinct points, and let \(W_d(Z)\) denote the object called the \emph{\(d\)-Weddle locus} in the supplied source. This notation is a placeholder for the source's construction: the supplied record does not define that construction, specify the ground field or its characteristic, or state the admissible ranges of \(n,d,r\). The word ``general'' is retained as a source hypothesis; a precise parameter space and genericity convention require confirmation.

Subject to those unresolved specifications, the requested problem is:
\[
\text{Determine }\dim W_d(Z)\text{ for a general set }Z\subset\mathbb P^n\text{ of }r\text{ points, as a function of }n,d,r.
\]
The supplied evidence does not suffice to replace \(W_d(Z)\) by defining equations or a rank condition. It also does not specify how an empty locus should be reported.

The original excerpt emphasizes an unresolved case in \(\mathbb P^3\), but its point count is rendered as ``\(d+2\;2-1\)''. This is corrupted mathematical typography; the intended cardinality cannot be asserted from the supplied text alone.

The excerpt additionally presents a related, separately numbered Question 2: which quartic surfaces in \(\mathbb P^3\) occur as the Weddle surface of a linear system \(S\) of quadrics of dimension \(3\)? In particular, does a general determinantal quartic surface arise from a suitable such system? The dimension convention for \(S\), the intended meaning of ``determinantal quartic'', and the construction of its Weddle surface are not supplied. These questions are contextual and are not identified with the principal dimension problem.

For context, the excerpt describes the Weddle surface determined by six points as special, mentioning its fifteen lines joining pairs of base points and referring to an unavailable reference [2]. It also says that the interpolation matrix for a general linear system of quadrics need not be symmetric, although it arises from a partially symmetric tensor. It proposes extensions to higher degrees and projective dimensions, with the interpolation matrix obtained by \(d-1\) tensor contractions. These source remarks do not provide a complete definition of the loci in the principal question.
\subsection{Short English statement}
Determine the dimension of the \(d\)-Weddle locus associated with \(r\) general points in \(\mathbb P^n\). A complete formulation requires the missing definition of that locus, ground-field assumptions, and parameter ranges; the highlighted three-dimensional special case also has a corrupted point count. The supplied partially solved classification is dated curation, not an independently verified status.
\subsection{Sources}
\begin{itemize}
\item Supplied local UnsolvedMath registry record, permanent ID \(30005298\), source code \texttt{catalogue entry}; registry collection \texttt{data/unsolvedmath/problems.json}, with supplementary existing records maintained in \texttt{catalog/problems.json}. Title: \emph{Dimensions of General Weddle Loci}. No individual proposer is named in the record.
\item Supplied report citation: \emph{Mini-Workshop: Subvarieties in Projective Spaces and Their Projections} (2023), Oberwolfach Reports, pp.\ 3093--3119, DOI \texttt{10.4171/owr/2022/54}. The extracted text contains the page marker 3103. Supplied source and literature-source link, labeled ``EMS Press, OWR 2022/54'': \url{https://doi.org/10.4171/owr/2022/54}. The record separately gives proposed year 2022; the discrepancy with the citation's year has not been resolved.
\item Supplied literature source: Chiantini et al., \emph{Weddle schemes}, \url{https://arxiv.org/abs/2606.25060}. The record supplies no fuller author list or theorem references.
\item Literature assessment dated exactly \texttt{2026-08-22}: the supplied curation labels the problem partially solved and reports that the 2026 Weddle-schemes paper determines the hypersurface case and its degree for general sets of binomially many points, studies nearby cardinalities, and leaves the dimension for arbitrary \(r\) open, including the cited unresolved \(\mathbb P^3\) case. This is supplied dated curation, not independently verified literature assessment; the binomial cardinality and precise theorem hypotheses are not supplied.
\item The registry labels the statement \texttt{corrected\_verified} and reports that it was checked against the official Oberwolfach report. That verification is a supplied registry claim, not verification performed for this draft. The excerpt's reference [2] has no corresponding bibliographic entry in the supplied record.
\end{itemize}
\end{document}
